% ── 消融研究（物理建模的贡献） ──
\subsection{物理建模的消融研究}
\label{sec:ablation}

为了量化入射角（$\theta$）与斜视角（${\psi_{\text{sq}}}$）建模对求解器性能的贡献，
我们使用四种目标函数变体重复了 MOEA-3 实验：
%
\begin{description}
\item[A. 完整物理模型（基线）] $f_2 = \frac{1}{|\mathcal{S}|}\sum \sin\theta_{\text{elev},i} \cdot \cos\psi_{\text{sq},i}$，
      $f_3 = \frac{1}{|\mathcal{S}|}\sum \cos^3\xi_i$
      （与 \S~\ref{sec:pareto-coupling} 中的标准配置相同）。
\item[B. 去除斜视分量] $f_2 = \frac{1}{|\mathcal{S}|}\sum \sin\theta_{\text{elev},i}$，
      $f_3 = \frac{1}{|\mathcal{S}|}\sum \cos^3\theta_{\text{elev},i}$（移除 ${\psi_{\text{sq}}}$ 分量；
      $\theta_{\text{elev}}$ 为高程面入射角，$\cos\theta_{\text{elev}}=\cos\xi/\cos\psi_{\text{sq}}$）。
\item[C. 去除入射角分量] $f_2 = \frac{1}{|\mathcal{S}|}\sum \cos\psi_{\text{sq},i}$，
      $f_3 = \frac{1}{|\mathcal{S}|}\sum \cos^3\psi_{\text{sq},i}$（移除 $\theta_{\text{elev}}$ 分量）。
\item[D. 不含物理信息] $f_2 = f_3 = 1$（逐任务的目标为常数；求解器不接收
      任何几何信息；表~\ref{tab:ablation} 中所示的 $f_2$、$f_3$
      为基于所得调度结果实际观测几何事后计算的诊断值）。
      变体 D 充当零模型对照：它检验 NSGA-III 求解器的性能应归因于 $f_2$ 和 $f_3$
      的物理内容，还是仅仅归因于结构化搜索的附加目标的存在。
\end{description}
%
四种变体均使用相同的 NSGA-III 参数（种群大小 $= 100$，
代数 $= 200$，Das--Dennis $p=12$）、相同的场景文件，
以及相同的 G-BL 热启动注入。仅目标函数不同，
因此本实验属于受控消融，而非重新调参实验。
在相同 NSGA-III 参数与相同 G-BL 种子注入下的求解器内消融充分控制了混淆因素，
使得 Wilcoxon 符号秩检验适用于目标函数变体的配对比较。
每种变体完成 6 个场景类别 $\times$ 50 个随机种子 $= 300$ 次运行
（A--D 共计 1,200 次运行）；表~\ref{tab:ablation} 报告密度区间类别 S1--S4。
其余两个类别施压于非密度轴：S5 保持 $N=20$ 密度但将目标位置沿轨散布于
$5^\circ$/$15^\circ$/$25^\circ$/$35^\circ$/$45^\circ$（斜视散布灵敏度组），
S6 将 20 个目标打包为 2--10 个窄窗簇（窗口宽 60--180 秒；C3 转换应力组）。
退化指标定义为
$\Delta\% = (\text{A} - \text{V})/\text{A} \times 100$，
其中 A 为完整物理基线，V 为待比较变体；
负值表示变体在 $f_1^*$ 上优于基线。

\begin{table}[H]
\centering
\caption{各场景类别下的消融结果。加粗：$p < 0.005$（与 A 配对 Wilcoxon 符号秩检验；与 \S\ref{sec:evaluation-metrics} 的 Friedman--Wilcoxon 框架一致）。}
\label{tab:ablation}
\begin{tabular}{llccccrc}
\toprule
\textbf{类别} & \textbf{V} & \textbf{$f_1^*$} & \textbf{$f_2$} &
\textbf{$f_3$} & \textbf{$f_1/n$} & \textbf{$n_\text{sel}$} &
\textbf{$\Delta f_1^*\%$} \\
\midrule
\multirow{4}{*}{S1 ($N=20$)}
 & A & $0.850 \pm 0.123$ & $0.376$ & $0.698$ & $5.99 \pm 0.73$ &
   $15.6 \pm 2.6$ & (基线) \\
 & B & $0.852 \pm 0.152$ & $0.353$ & $0.574$ & $6.01 \pm 0.83$ &
   $15.7 \pm 3.1$ & $-0.8\%$ \\
 & C & $0.845 \pm 0.095$ & $0.375$ & $0.703$ & $5.86 \pm 0.80$ &
   $16.0 \pm 2.2$ & $+0.1\%$ \\
 & D & $1.002 \pm 0.011$ & $0.302$ & $0.525$ & $5.65 \pm 0.66$ &
   $19.5 \pm 0.7$ & $\mathbf{-18.0\%}$ \\
\midrule
\multirow{4}{*}{S2 ($N=100$)}
 & A & $0.976 \pm 0.040$ & $0.329$ & $0.597$ & $5.58 \pm 0.42$ &
   $62.7 \pm 19.3$ & (基线) \\
 & B & $0.990 \pm 0.022$ & $0.321$ & $0.570$ & $5.55 \pm 0.37$ &
   $63.6 \pm 19.1$ & $\mathbf{-1.2\%}$ \\
 & C & $0.977 \pm 0.033$ & $0.328$ & $0.599$ & $5.57 \pm 0.40$ &
   $62.8 \pm 19.3$ & $+0.0\%$ \\
 & D & $1.000 \pm 0.000$ & $0.317$ & $0.562$ & $5.54 \pm 0.37$ &
   $64.2 \pm 18.8$ & $\mathbf{-2.0\%}$ \\
\midrule
\multirow{4}{*}{S3 ($N=300$)}
 & A & $0.998 \pm 0.005$ & $0.344$ & $0.623$ & $5.63 \pm 0.30$ &
   $154.2 \pm 39.9$ & (基线) \\
 & B & $0.998 \pm 0.005$ & $0.343$ & $0.618$ & $5.63 \pm 0.29$ &
   $154.2 \pm 40.0$ & $+0.0\%$ \\
 & C & $0.998 \pm 0.004$ & $0.344$ & $0.623$ & $5.63 \pm 0.30$ &
   $154.2 \pm 39.9$ & $+0.0\%$ \\
 & D & $1.000 \pm 0.000$ & $0.342$ & $0.616$ & $5.63 \pm 0.29$ &
   $154.5 \pm 39.9$ & $\mathbf{-0.2\%}$ \\
\midrule
\multirow{4}{*}{S4 ($N=500$)}
 & A & $0.997 \pm 0.008$ & $0.338$ & $0.680$ & $5.83 \pm 0.44$ &
   $142.3 \pm 72.1$ & (基线) \\
 & B & $0.999 \pm 0.003$ & $0.337$ & $0.677$ & $5.84 \pm 0.44$ &
   $142.5 \pm 72.0$ & $-0.1\%$ \\
 & C & $0.998 \pm 0.006$ & $0.337$ & $0.680$ & $5.83 \pm 0.44$ &
   $142.4 \pm 72.0$ & $+0.0\%$ \\
 & D & $1.000 \pm 0.000$ & $0.336$ & $0.675$ & $5.83 \pm 0.45$ &
   $142.7 \pm 71.9$ & $\mathbf{-0.2\%}$ \\
\bottomrule
\end{tabular}
\end{table}

\begin{figure}[H]
    \centering
    \includegraphics[width=\textwidth]{fig5_ablation.pdf}
    \caption{四种目标函数变体（B--D）与四类场景（S1--S4）下的消融结果。柱状图表示相对完整物理基线 A 的归一化收益 $f_1^*$ 与事后几何质量指标 $f_2$、$f_3$ 的退化（负值表示变体优于 A）；星号标记与 A 配对 Wilcoxon 比较 $p<0.005$（Bonferroni 阈值）。}
    \label{fig:fig5}
\end{figure}

表~\ref{tab:ablation} 与图~\ref{fig:fig5} 中的消融结果表明，物理目标建模的影响
集中在稀疏场景中，并随场景密度的增大而消失 ---
这与我们的尺度敏感性分析（\S~\ref{sec:scale-sensitivity}）模式一致。

在 S1（$N = 20$）下，仅变体 D（不含物理信息）在本文 $p<0.005$ 阈值下
与基线 A 存在显著差异：其 $f_1^*$ 升至 $1.002$（$-18.0\%$，$p=3.4\times10^{-9}$），
通过多调度 $3.9$ 个任务（$19.5$ vs.\ $15.6$），代价是事后 $f_3$ 从 $0.698$
降至 $0.525$（$-25\%$）。变体 B（去除斜视）使 $f_1^*$ 基本不变（$-0.8\%$，$p=0.46$），
但将 $f_3$ 降至 $0.574$（$-18\%$）、$f_2$ 降至 $0.353$（$-6.2\%$，$p=1.1\times10^{-5}$）：
缺少斜视惩罚后求解器容忍更高的斜视视角，牺牲了完整物理基线在稀疏区域购买的辐射质量。
变体 C（去除入射角分量）在报告的任何指标上均无可测变化。
这种惰性与窗口几何一致：高程面入射角在一个可见性窗口内仅变化几度（S1 场景中均值约 $1.5^\circ$），
而斜视角扫过约 $20$--$35^\circ$，因此 $f_2$/$f_3$ 中的 $\theta_{\text{elev}}$ 因子在每个窗口内近似为共同尺度；
移除它保持了窗口内任务排序，从而保持了选中的调度。
\S\ref{sec:pareto-coupling} 的入射角冲突因此通过跨窗口/聚合几何起作用，而非通过窗口内选择。
因此，物理目标的内容在稀疏区主要通过辐射质量（$f_3$）与完全移除时的覆盖率反差影响结果，
而仅移除任一几何分量均不改变覆盖率。

到 S2（$N = 100$）时，变体 B 与 D 仍与 A 显著不同但效应很小
（$-1.2\%$，$p=5.3\times10^{-4}$ 与 $-2.0\%$，$p=1.7\times10^{-6}$）：
移除斜视或全部物理信息仍会适度改变覆盖率。变体 C 仍与 A 无法区分。

在 S3--S4（$N \geq 300$）下，每种变体在 $f_1^*$ 上均落在基线 $0.2\%$ 以内，
在 $f_2$ 上均在 $0.6\%$ 以内；D 的配对检验显著性反映的是一致但可忽略的效应（$-0.2\%$）。
这种数值不变性正是 \S~\ref{sec:scale-sensitivity} 中讨论的
尺度依赖饱和的直接表征：一旦候选观测窗口的密度饱和了调度，
目标函数便仅惩罚不可行的解，而不再区分可行的解。
在此情形下，无论 $f_2$ 与 $f_3$ 是否承载几何含义（A），
还是充当常数占位符（D），NSGA-III 求解器
本质上等价于一种约束过滤下的收益最大化器。

两个稳健性类别在密度轴之外复制了该模式。在 S5（沿轨目标散布）和 S6（窄窗簇）上，
完全移除变体 D 在 $f_2$ 上显著差于 A（$-11.2\%$，$p=1.6\times10^{-13}$ 和 $-11.1\%$，$p=1.8\times10^{-15}$）、
在 $f_3$ 上同样更差（$-6.1\%$ 和 $-6.3\%$），C 与 A 不可区分（$f_2$ $+0.4\%$，$p=0.33$ 和 $-0.6\%$，$p=0.48$），
B 呈现小幅显著差异（$f_2$ $+3.2\%$，$p=0.046$ 和 $+4.6\%$，$p=1.2\times10^{-5}$）。
物理目标的覆盖率代价在此同样持续：D 比 A 多获得 $17\%$（S5）和 $8\%$（S6）的原始收益，
即在稀疏配置下目标驱动的选择无论施压轴是密度、目标散布还是窗口簇都会牺牲覆盖率。

综合来看，物理质量目标主要在稀疏区域增加价值，且主要通过辐射质量：
在 S1，完整物理变体 A 相比无物理变体 D 实现 $f_3$ 的 $33\%$ 提升（$0.698$ vs.\ $0.525$；在共同退化基线上 D 比 A 低 $25\%$），
代价是 $18\%$ 覆盖率劣势；仅移除斜视分量（B）已在保持 $f_1^*$ 的同时牺牲 $18\%$ 的 $f_3$。
部分移除目标的覆盖率效应在所有区间均极小，且到 $N \geq 300$ 时
物理目标的选择对结果数值上不再有影响。
